%%%%%%%%%%%%%%%%%%%%%%%%%%%%%%%%%%%%%%%%%
% FRI Data Science_report LaTeX Template
% Version 1.0 (28/1/2020)
%
% Jure Demšar (jure.demsar@fri.uni-lj.si)
%
% Based on MicromouseSymp article template by:
% Mathias Legrand (legrand.mathias@gmail.com)
% With extensive modifications by:
% Antonio Valente (antonio.luis.valente@gmail.com)
%
% License:
% CC BY-NC-SA 3.0 (http://creativecommons.org/licenses/by-nc-sa/3.0/)
%
%%%%%%%%%%%%%%%%%%%%%%%%%%%%%%%%%%%%%%%%%

%----------------------------------------------------------------------------------------
%  PACKAGES AND OTHER DOCUMENT CONFIGURATIONS
%----------------------------------------------------------------------------------------
\documentclass[fleqn,moreauthors,10pt]{ds_report}
\usepackage[english]{babel}
\usepackage{titlesec}
\usepackage{comment}
\usepackage{float}
\titlespacing*{\subsection}{0pt}{2.5ex plus 3pt minus 2pt}{0.25em}
\usepackage{svg}

\graphicspath{{fig/}}

%----------------------------------------------------------------------------------------
%  ARTICLE INFORMATION
%----------------------------------------------------------------------------------------

% Header
\JournalInfo{FRI Natural language processing course 2026}

% Interim or final report
\Archive{Project report}
%\Archive{Final report}

% Article title
\PaperTitle{A RAG-Based Chatbot for Slovenian Tax and Investment Law}

% Authors (student competitors) and their info
\Authors{Maj Bregar, Edita Džubur and Jan Tršan }

% Advisors
\affiliation{\textit{Advisors: Slavko Žitnik}}

% Keywords
\Keywords{LLM, retrieval-augmented generation, chatbot, prompt engineering}
\newcommand{\keywordname}{Keywords}

%----------------------------------------------------------------------------------------
%  ABSTRACT
%----------------------------------------------------------------------------------------

\Abstract{
  This project develops a domain-specific chatbot for Slovenian investment and tax law. Since general-purpose LLMs lack reliable domain knowledge, the proposed approach uses retrieval-augmented generation (RAG) to produce accurate, up-to-date, and verifiable responses based on external legal documents. Prompt engineering is applied to enforce structure and introduce guardrails. The system is built on a curated dataset combining legal corpora and official Slovenian sources.
}

%----------------------------------------------------------------------------------------

\begin{document}

% Makes all text pages the same height
\flushbottom

% Print the title and abstract box
\maketitle

% Removes page numbering from the first page
\thispagestyle{empty}

%----------------------------------------------------------------------------------------
%  ARTICLE CONTENTS
%----------------------------------------------------------------------------------------

\section*{Introduction}

Large language models (LLMs) often struggle in specialized domains where answers require precise, up-to-date, and verifiable information. This is particularly challenging in legal and financial settings, where responses may depend on specific regulations, exceptions, and recent legislative changes. In such applications, generated answers should not only be plausible, but also grounded in inspectable sources.

Retrieval-Augmented Generation (RAG), introduced by Lewis et al.~\cite{rag}, combines language generation with external document retrieval and has become increasingly attractive for knowledge-intensive applications. Recent work has explored retrieval-based systems for legal question answering and consultation. LegalBench-RAG~\cite{legalbenchrag} highlighted the importance of accurate retrieval in legal systems and showed that retrieval quality strongly affects downstream answer generation. LexRAG~\cite{lexrag} further studied multi-turn legal consultation and identified retrieval as a major challenge in conversational legal assistants. Domain-specific legal assistants have also been proposed for other jurisdictions, such as the Indian system Grahak-Nyay~\cite{grahaknyay}, a consumer-law chatbot based on large language models.

Motivated by these findings, we develop a RAG-based chatbot for Slovenian investment and tax law. The system retrieves relevant legal documents and generates grounded responses based on Slovenian legal sources.



%----------------------------------------------------------------------------------------
%  METODE
%----------------------------------------------------------------------------------------


\newpage
\section*{Methods}

\subsection*{Dataset}

The retrieval corpus is based on the Slovenian legal dataset released within the LLM4DH project~\cite{llm4dh}. We use the COLESLAW corpus and specifically the collection \emph{register-predpisov}, which contains structured Slovenian legal acts and regulatory documents.

The corpus was filtered to retain tax- and investment-related legislation, including laws such as \textit{ZDoh-2} and \textit{ZDavP-2}. Documents were converted into a unified textual representation before preprocessing and chunk generation.



\subsection*{System Architecture}

The chatbot follows a modular Retrieval-Augmented Generation (RAG) architecture consisting of offline indexing and online inference stages. During indexing, legal documents are processed and embedded into a vector store. During inference, relevant chunks are retrieved, reranked, and passed to a language model for grounded response generation.

Figure~\ref{fig:pipeline} illustrates the overall pipeline.

%-------------------------
% PLACEHOLDER FIGURE
%-------------------------

\begin{figure}[h]
  \centering
  \fbox{
    \parbox{0.85\linewidth}{
      \centering

      Legal documents

      $\downarrow$

      Chunking + metadata extraction

      $\downarrow$

      Embedding (BGE-M3)

      $\downarrow$

      Cosine Similarity Retrieval (top 50 chunks)

      $\downarrow$

      Reranking and final chunk output (BGE-reranker-v2-M3, top 5 chunks)

      $\downarrow$

      Prompt construction

      $\downarrow$

      Qwen2.5-14B-Instruct LLM response generation
  }}
  \caption{Overview of the RAG pipeline.}
  \label{fig:pipeline}
\end{figure}

\subsection*{Retrieval and Generation Pipeline}

Legal documents were segmented according to their internal structure into article- and paragraph-level chunks. Metadata including law title, article number, and legal categories was attached to each chunk to preserve source context.

Embeddings were generated using the \emph{BAAI/bge-m3}~\cite{bgem3} multilingual embedding model. During inference, retrieved candidates were ranked by cosine similarity before being reranked by the \emph{BAAI/bge-reranker-v2-m3}~\cite{reranker} and passed to the generation stage. Multiple reranking strategies were explored and are discussed in Results.

%-----------------------------------

The final prompt included the top retrieved chunks together with metadata and conversation history. Prompt engineering was used to improve grounding behavior. Instead of relying exclusively on the highest-ranked retrieved chunk, the prompt instructed the model to consider all retrieved evidence jointly and resolve conflicting information before generating a response. We additionally experimented with different prompt formulations and chunk-selection strategies.

Responses were generated using \emph{Llama-3.1-8B-Instruct}~\cite{llmllama} and \emph{Qwen2.5-14B-Instruct}~\cite{qwen}, served through vLLM.


\subsection*{Retriever Evaluation Setup}

To select the final reranking model, we conducted an LLM-supervised retrieval quality evaluation using the NDCG@k metric. We constructed a set of handwritten queries representing questions about Slovene tax and investment legislation. For each query, we first retrieved 100 chunks using cosine similarity ranking on embeddings generated by the \textbf{BGE-M3} embedding model. The retrieved chunks were then evaluated using \textit{gpt-4.1-mini}, which was prompted to assess how relevant each chunk was for answering the query, as well as how effectively the chunk could support answer generation. The resulting LLM-generated relevance scores were treated as the ground-truth ranking for computing the NDCG metric on the final retrieved chunks.

We first evaluated the baseline retrieval performance of the \textbf{BGE-M3} embedding model without reranking. We then evaluated retrieval performance after applying reranking with three reranking models of different sizes, shown in Table~\ref{tab:reranker_models}.

\begin{table}[h]
\centering
\caption{Reranking model names}
\label{tab:reranker_models}
\renewcommand{\arraystretch}{1.2}
\small
\begin{tabular}{|l|}
\hline
\textbf{Reranker Model} \\
\hline
\textbf{cross-encoder/ms-marco-MiniLM-L6-v2} \\
\hline
\textbf{cross-encoder/mmarco-mMiniLMv2-L12-H384-v1} \\
\hline
\textbf{BAAI/bge-reranker-v2-m3} \\
\hline
\end{tabular}
\end{table}

Each model configuration was evaluated using the mean NDCG@k value across 30 unique query results.





\subsection*{LLM Answer Evaluation Setup}

To evaluate answer quality, we created a gold dataset of 29 Slovenian tax-related questions. The dataset includes realistic user scenarios, definition questions, edge cases, and a small number of unanswerable questions. It is designed to test both retrieval quality and whether the model can refrain from answering when there is insufficient legal evidence. The question--answer pairs were first generated using LLMs and then manually verified against Slovenian legal sources.

The overall evaluation process is summarized in Figure~\ref{fig:evaluation_pipeline}.

\begin{figure}[h]
  \centering
  \fbox{
    \parbox{0.85\linewidth}{
      \centering

      Gold dataset (29 Slovenian tax questions)

      $\downarrow$

      LLM-generated Q\&A pairs

      $\downarrow$

      Manual verification against legal sources

      $\downarrow$

      Evaluation inputs:
      (gold answer + model answer + retrieved context)

      $\downarrow$

      LLM judge (GPT-5-nano)

      $\downarrow$

      JSON output:
      \textbf{legal\_correct}, \textbf{citation\_correct}, \textbf{hallucination}, \textbf{passed}, \textbf{rationale}

      $\downarrow$

      Manual review (if needed)
    }
  }
  \caption{Overview of the LLM-based evaluation pipeline.}
  \label{fig:evaluation_pipeline}
\end{figure}

We evaluated system outputs using an automatic LLM-based judge (\emph{GPT-5-nano}). For each question, the judge receives the gold answer, generated answer, and retrieved legal context. It produces a structured JSON output with five fields: \textbf{legal\_correct}, \textbf{citation\_correct}, \textbf{hallucination}, \textbf{passed}, and \textbf{rationale}, which together provide both a binary decision and an explanation of the judgment.

Initial results showed that the automated judge was sometimes overly strict, penalizing answers that were practically correct but included additional citations or standard legal disclaimers. To address this limitation, we added a manual review step. Manual evaluation focuses on whether a user relying only on the answer would reach the correct legal conclusion.


%----------------------------------------------------------------------------------------
%  REZULTATI
%----------------------------------------------------------------------------------------

\newpage
\section*{Results}

\subsection*{Retriever Evaluation Results}

The retrieval evaluation results are shown in Figure~\ref{fig:reranking_means}. The results demonstrate that reranking has a substantial impact on retrieval quality. The largest and most recent reranking model, \textit{BAAI/bge-reranker-v2-m3}, consistently achieved the highest NDCG@k scores across all evaluated values of $k$, significantly outperforming both the baseline embedding-only retrieval and the alternative reranking models.

In contrast, the smaller \textit{cross-encoder/ms-marco-MiniLM-L6-v2} model reduced retrieval performance compared to the baseline. This suggests that low-capacity reranking models can negatively affect retrieval quality when paired with semantically strong embedding models such as \textit{BGE-M3}, as incorrect reranking disrupts the high quality original ordering.

The NDCG@k trends also provide insight into the optimal number of retrieved chunks for the final pipeline. The highest scores are achieved for small values of $k$, particularly between $k=2$ and $k=5$, indicating that the most semantically relevant chunks are typically ranked within the top retrieved results. Beyond this range, the scores gradually increase again due to the inclusion of weakly relevant chunks, which artificially inflate the ranking metric without meaningfully improving retrieval quality.

Although NDCG@k is a relative ranking metric, the observed scores remain substantially below the theoretical maximum value of 1.0. This indicates that the retrieval system does not consistently return perfectly relevant chunks at the top ranks. This observation is further supported by the retrieval variance shown in Figure~\ref{fig:reranking_variance}, where large score variations at low values of $k$ reveal inconsistent retrieval quality across evaluation queries.

\begin{figure}[H]
  \centering
  \includesvg[width=\linewidth]{fig/model_reranker_llm_scored_mean}
  \caption{
    Retrieval performance of embedding and reranking model combinations measured using average NDCG@k over 30 evaluation queries.
  }
  \label{fig:reranking_means}
\end{figure}

\begin{figure}[H]
  \centering
  \includesvg[width=\linewidth]{fig/model_reranker_llm_scored_std}
  \caption{
    Standard deviation of NDCG@k scores across evaluation queries for different embedding and reranking model combinations.
  }
  \label{fig:reranking_variance}
\end{figure}

Based on both the mean and variance of the NDCG@k scores, the \textit{BAAI/bge-reranker-v2-m3} model was selected for the final pipeline. To reduce retrieval latency, we set the cosine similarity retrieved number of chunks to $50$ in the final pipeline, under the assumption that most high quality embeddings are captured within that window. We additionally selected $k=5$ as the final number of retrieved chunks, allowing the downstream LLM to identify and utilize the most relevant information during answer generation. 




\subsection*{Question Answering Evaluation Results}

We evaluated Llama-3.1 and Qwen-2.5 across 29 test questions, comparing strict automated decisions with our practical manual review. The final metrics are summarized in Table~\ref{tab:qa_comparison}.

\begin{table}[h]
\centering
\small
\setlength{\tabcolsep}{6pt}
\renewcommand{\arraystretch}{1.1}
\begin{tabular}{lcccc}
\hline
 & \multicolumn{2}{c}{\textbf{Llama-3.1 (8B)}} & \multicolumn{2}{c}{\textbf{Qwen-2.5}} \\
\textbf{Metric} & Auto & Manual & Auto & Manual \\
\hline
Passed (True)   & 12    & 17    & 20   & 22     \\
Failed (False)  & 17   & 12      & 9   & 7      \\
\textbf{Pass Rate} & \textbf{41\%} & \textbf{58\%} & \textbf{69\%} & \textbf{76\%} \\
\hline
\end{tabular}
\caption{Comparison of true pass results for Llama and Qwen models.}
\label{tab:qa_comparison}
\end{table}

The automated judge frequently penalized both models for harmless additions such as extra legal references or standard legal disclaimers. Manual review showed substantially stronger practical performance. Qwen-2.5 achieved the highest overall performance, correctly answering 22 of 29 questions, while Llama-3.1 correctly answered 17. The gap between automated and manual scores suggests that strict citation-based evaluation may underestimate practical answer quality.

Manual inspection also uncovered errors that were not captured by the automated judge. In one example, the model was asked whether excise duty (slo. \emph{trošarina}) must be paid for 8 liters of spirits (slo. \emph{žgane pijače}) brought from Italy for personal use. The model correctly identifies the 10-liter threshold for excise duty on spirits but incorrectly applies this rule, concluding that duty is required even though the quantity is below the legal limit. This shows that errors can occur even when the relevant legal rule is correctly retrieved, due to incorrect application of the rule to the specific scenario.
%----------------------------------------------------------------------------------------
%  DISCUSI
%----------------------------------------------------------------------------------------

\section*{Discussion}

The results indicate that retrieval quality and language model choice both significantly affect performance in legal question answering systems. Our analysis shows that the quality of the final answer is only as good at the weakest link in the entire pipeline, therefore a selected strong baseline embedding model must be paired with a strong reranking model and a powerful large language model.

While reranking improved the relevance of retrieved context, this improvement did not always translate into better final answers. In several cases, the model still failed to correctly use relevant legal passages, suggesting that improving retrieval alone is not sufficient. The downstream model must be able to reliably interpret and select the relevant information from multiple sources.

We also observed clear differences between models. Qwen-2.5 generally produced more accurate and consistent answers than Llama-3.1 in practical evaluations. However, even the stronger model was not robust in all cases, particularly for edge cases or questions requiring careful interpretation of legal conditions. This shows that model choice reduces but does not eliminate reasoning and grounding errors.

Prompt design had a strong influence on system behaviour, but introduced trade-offs. Early prompts caused excessive reliance on the top-ranked retrieved chunk, leading to incomplete answers. After modifying the prompt to encourage the use of multiple sources, the system became more comprehensive but also less precise. In some cases, it combined unrelated information from different documents, which resulted in overly broad or partially incorrect legal conclusions.

A major limitation of the system is the instability of smaller language models when operating with retrieved context. They occasionally misinterpreted simple user queries, introduced irrelevant explanations, or produced logically inconsistent reasoning despite having access to correct legal sources. This suggests that retrieval-augmented generation alone is not sufficient for reliable legal reasoning without stronger instruction-following or reasoning capabilities.

Another limitation concerns evaluation. The automated LLM-based judge (\emph{GPT-5}) was not fully aligned with human judgment. It tended to penalise answers that included additional legal disclaimers or extra citations, even when the practical answer was correct. This required manual verification, which indicates that automated evaluation is not fully reliable for this type of task.


Future evaluation improvements could include a more careful separation between retrieval quality and answer evaluation. Specifically, evaluation criteria should be defined more strictly to distinguish whether the system arrives at the right legal conclusion or whether it's citation is correct. 

Future pipeline improvements could include expanding the dataset with more complex legal scenarios, integrating expert legal review for evaluation, and testing stronger reasoning-oriented models. Another promising direction would be improving retrieval beyond semantic similarity, for example by incorporating structured legal metadata or rule-based filtering to reduce irrelevant context.


\bibliographystyle{unsrt}
\bibliography{report}

\end{document}
